\section{Governança de TI}

\recado{Citar fonte mais diretamente}
\revisando{
A governança de TI pode ser compreendida como o conjunto de estruturas, processos e mecanismos 
institucionais estruturados para definir direitos decisórios, responsabilidades e formas de 
controle, visando garantir o pleno alinhamento do uso da tecnologia aos objetivos organizacionais \cite{governanca_ia_2026}.
}

\revisado{
Nessa perspectiva, ela estabelece a especificação dos direitos de decisão e o arcabouço de 
responsabilidades necessários para incentivar comportamentos desejáveis no uso da tecnologia. Assim, 
a governança de TI reflete os princípios amplos da governança corporativa, focando na administração 
e no direcionamento da TI para o alcance das metas de desempenho da organização \cite{parcial_2025_processo_gov_ti}.
}

\revisado{
Sob a ótica da alta administração, o termo descreve a forma como os líderes e responsáveis pela 
governança organizacional consideram a TI nas funções de supervisão, monitoramento, controle e 
direcionamento. A maneira como a tecnologia é aplicada impacta diretamente a capacidade da 
organização de atingir sua visão, missão e objetivos estratégicos \cite{rondonia_2023}.
}

\revisado{
Por fim, é essencial diferenciar esse conceito do nível operacional: enquanto a gestão de TI foca na 
execução e na administração cotidiana dos recursos tecnológicos, a governança de TI atua no plano 
estratégico, sendo responsável por definir prioridades, pactuar responsabilidades e instituir 
mecanismos de controle e monitoramento das decisões de tecnologia \cite{governanca_ia_2026}.
}

\revisado{
A governança de TI originou-se da governança corporativa e é composta por estruturas, processos 
organizacionais e elementos de liderança que asseguram a sustentação e a expansão das estratégias e 
metas organizacionais por meio da tecnologia. Por essa razão, prevalece na literatura o entendimento 
de que a governança de TI constitui um subconjunto da governança corporativa \cite{rangel_2021}.
}

\revisado{
O fortalecimento dessa relação foi impulsionado pelo surgimento de normativas rígidas, a exemplo do 
Ato Sarbanes-Oxley, promulgado nos Estados Unidos em 2002, que forçou diversas organizações a 
reestruturarem suas práticas de governança corporativa para assegurar controles fiscais adequados 
aos acionistas e demais partes interessadas \cite{rangel_2021}.
}

\revisado{
Nesse contexto, a governança viabiliza o alinhamento entre a TI e o negócio, garantindo que as 
deliberações da alta direção considerem aspectos como gestão de riscos, planejamento financeiro, 
continuidade dos negócios e conformidade regulatória. Com isso, a TI deixa de ser percebida como um 
mero centro de custos ou instrumento de suporte e passa a integrar a gestão estratégica e financeira 
da organização \cite{joyce_mariana_pereira_2023_7555947}.
}

\revisado{
A governança é responsável por assegurar que as necessidades, condições e opções das partes 
interessadas sejam devidamente avaliadas para estabelecer objetivos corporativos equilibrados e 
pactuados. A partir disso, define-se o direcionamento estratégico mediante priorizações e tomadas de 
decisão, com o subsequente monitoramento do desempenho e da conformidade frente aos alvos fixados 
\cite{joyce_mariana_pereira_2023_7555947}.
}

\revisado{
Para atingir esses objetivos, a Governança de TI estrutura-se em áreas de foco interconectadas, 
como alinhamento estratégico, entrega de valor, gestão de riscos, gestão de recursos e mensuração de 
desempenho, sendo a aplicação conjunta dessas etapas indispensável para a obtenção dos melhores 
resultados organizacionais \cite{rondonia_2023}.
}

\revisado{
Nesse contexto, destacam-se dimensões fundamentais para a consolidação do papel da tecnologia no 
negócio. O alinhamento estratégico assegura que as ações de TI promovam a missão organizacional, 
enquanto a entrega de valor garante a execução de projetos dentro do prazo e do orçamento previstos. 
Paralelamente, a gestão de recursos otimiza os investimentos e os ativos de informação, a gestão de 
riscos protege o patrimônio tecnológico e informacional da instituição, e a mensuração de desempenho 
monitora continuamente a implementação das estratégias para maximizar o valor gerado pela TI \cite{rangel_2021}.
}

\revisado{
Para que os objetivos e áreas de foco da tecnologia sejam efetivamente operacionalizados, a 
literatura consolida que os mecanismos de governança de TI se organizam em três dimensões 
complementares: estruturas, processos e relacionamentos. As estruturas englobam os comitês de 
governança digital, unidades de TI, escritórios de projetos, áreas de segurança e fóruns decisórios. 
Os processos referem-se ao planejamento, priorização de demandas, gestão de portfólio, riscos e 
benefícios, além do acompanhamento de indicadores e auditoria. Por fim, os mecanismos relacionais 
conectam a TI às áreas finalísticas por meio de comunicação integrada, capacitação contínua, 
engajamento da alta administração e construção de uma linguagem comum na organização 
\cite{mello_ramos_2026_21402664}.
}

\subsection{Governança de TI na Administração Pública}

\revisado{
A aplicação da governança no contexto governamental exige compreender que as estratégias do setor 
público diferem substancialmente das do setor privado. Enquanto as empresas privadas buscam o 
retorno sobre o capital e o valor aos acionistas, as instituições públicas voltam-se à entrega de 
valor público. Dessa forma, os investimentos estatais em tecnologia devem ser avaliados com base na 
efetividade do atendimento à sociedade e no cumprimento das metas fixadas nos programas 
governamentais, extrapolando métricas puramente econômicas \cite{neto_avaliacao_2020}.
}

\revisado{
Embora compartilhem pilares conceituais básicos, os modelos de governança de TI em ambos os setores 
distinguem-se em razão do ambiente institucional. No setor público, prevalecem exigências de 
transparência, legitimidade das ações perante a sociedade e conformidade a dispositivos legais e 
órgãos de controle, que sobrepõem-se ao foco em resultados financeiros característico da iniciativa 
privada \cite{sena_governancti_2018}.
}

\revisado{
Essas particularidades ambientais manifestam-se em procedimentos burocráticos rígidos, tomadas de 
decisão altamente estruturadas, dependência de políticas públicas, elevada responsabilização da 
administração e rotatividade de lideranças em virtude de cargos temporários. Esse conjunto de 
contingências afeta diretamente a condução da tecnologia e demanda arranjos de governança 
ajustados à cultura do setor público \cite{sena_governancti_2018}.
}

\revisado{
Nesse cenário, a condução da governança de TI constitui responsabilidade indelegável da Alta 
Administração - compreendendo a liderança política e o corpo executivo. Cabe aos ocupantes do topo 
da organização fixar diretrizes decisórias, gerir riscos e coordenar estrategicamente as ações 
tecnológicas, evitando a falha recorrente de transferir o papel de governança estritamente para a 
área operacional de tecnologia \cite{sena_governancti_2018}.
}

\revisado{
Entretanto, as instituições públicas enfrentam entraves para efetivar essa gestão. Na administração 
pública, a escassez de recursos humanos e de profissionais com competências gerenciais dificulta a 
separação entre as atribuições estratégicas de governança e as rotinas operacionais de gestão. 
Além disso, o foco predominantemente técnico dos concursos públicos resulta na escassez de 
habilidades de planejamento e governança entre os servidores \cite{torriani_modelo_2019}.
}

\recado{Citar fonte mais diretamente}
\revisando{
Outro desafio reside no fato de que frameworks e padrões de boas práticas amplamente difundidos no 
mercado, como COBIT e ITIL, foram concebidos sob a ótica de organizações comerciais lucrativas, 
colidindo com as restrições orçamentárias e a atividade-fim governamental. Assim, a governança 
pública requer a simplificação e a adaptação dessas estruturas à realidade estatal para garantir a 
viabilidade das práticas \cite{sengik_governanctecnologia_2022}.
}

\revisado{
Em razão dessas barreiras, a presença de modelos formais e frameworks internacionais não assegura, 
por si só, uma governança efetiva no setor público. É comum observar a implementação parcial ou 
meramente documental de mecanismos de TI. Os obstáculos consolidados na literatura incluem escassez 
de recursos, baixo engajamento da alta administração, limitações na gestão e reduzida maturidade 
organizacional \cite{sena_governancti_2018}.
}

% { Removido -
% Para induzir a superação desses déficits na administração pública brasileira, o Tribunal de Contas 
% da União (TCU) instituiu levantamentos periódicos de avaliação e monitoramento. Essa medição 
% sistematizada de indicadores de desempenho visa mensurar a maturidade da governança de TI e orientar 
% as instituições públicas na busca por resultados e no aprimoramento da prestação de contas à 
% sociedade \cite{neto_avaliacao_2020}.
% }

%% colocar uma fala de conexão dos paragrafos

\revisar{
A avaliação e o monitoramento da governança pública pelo Tribunal de Contas da União tiveram início 
com levantamentos de perfil focados em Tecnologia da Informação, delineados com o propósito 
explícito de apoiar a melhoria contínua do desempenho das organizações públicas \cite{tcu_referencial_2020}. 
}

\revisar{
Realizados por meio de diagnósticos e questionários eletrônicos, esses levantamentos geram 
indicadores que permitem às instituições identificar vulnerabilidades em suas estruturas, subsidiar 
a definição de objetivos de aprimoramento e priorizar práticas de governança capazes de mitigar 
riscos que comprometam as entregas e a geração de valor para a sociedade \cite{tcu_referencial_2020}.
}

\revisado{
Estudos empíricos aplicados ao setor público nacional corroboram que as dimensões "Estratégias e 
Planos", "Liderança" e "Pessoas" são as variáveis de maior peso na diferenciação dos níveis de 
maturidade de governança de TI nas instituições. Fica evidente, portanto, que os órgãos públicos 
precisam focar esforços prioritários no alinhamento estratégico, na formação contínua e no 
desenvolvimento de competências de seus quadros para elevar seu grau de maturidade 
\cite{campelo_medeiros_maturidade_2016}.
}


\subsection{Mecanismos de Governança de TI}

\revisado{
A operacionalização da governança de TI nas organizações exige a implementação coordenada de um 
conjunto de mecanismos práticos, classificados pela literatura em três dimensões complementares: 
estruturais, processuais e relacionais. Enquanto as deliberações de alto nível fixam diretrizes e 
prioridades, cabe a esses mecanismos operacionalizar a execução cotidiana e o controle das decisões 
tecnológicas, permitindo a efetiva transição do plano estratégico para a prática administrativa 
\cite{governanca_ia_2026}.
}

\revisando{
No modelo de governança de TI proposto por \cite{sengik_governanctecnologia_2022} para a 
Administração Pública, esses mecanismos são organizados no Bloco de Operacionalização em três eixos 
de ação institucional.
}


\begin{itemize}
    \item \revisando{\textbf{Mecanismos Estruturais}: referem-se aos arranjos organizacionais formais e comitês 
    institucionais encarregados de conectar a alta administração, as áreas finalísticas e a área de 
    TI. Incluem a criação do Comitê Diretivo de TI, a especificação clara de papéis e 
    responsabilidades decisórias e a participação formal dos gestores de TI em instâncias superiores 
    de decisão.}

    \item \revisando{\textbf{Mecanismos Processuais}: envolvem a formalização, padronização e institucionalização das 
    rotinas de tomada de decisão e monitoramento das ações tecnológicas. Abrangem a definição de 
    indicadores de desempenho, o Planejamento Estratégico de TI, a formulação de Políticas de 
    Segurança da Informação, as práticas de gestão de riscos, a gestão do portfólio de projetos, a 
    priorização de investimentos e a gestão de serviços de TI.}
    
    \item \revisando{\textbf{Mecanismos Relacionais}: englobam os instrumentos formais e informais direcionados a 
    promover a comunicação, a liderança, o aprendizado compartilhado e a colaboração ativa entre os 
    atores de TI e os gestores de negócio. Contemplam a colaboração entre stakeholders, o 
    compartilhamento de conhecimento, a comunicação efetiva, o exercício da liderança de TI e as 
    parcerias estratégicas entre a TI e as áreas finalísticas.}
\end{itemize}


\revisado{
A eficácia do modelo depende da articulação harmônica entre esses três tipos de mecanismos e as 
áreas de foco da governança de TI, a exemplo do alinhamento estratégico, da gestão de riscos, da 
entrega de valor, do gerenciamento de recursos, da mensuração do desempenho e da prestação de contas. 
Nenhum mecanismo funciona de modo isolado: o desempenho de uma instância estrutural está 
condicionado ao suporte de processos bem delimitados e à fluidez dos canais relacionais que garantem 
o engajamento das partes interessadas \cite{sengik_governanctecnologia_2022}.
}

\revisado{
Com isso, a adoção da taxonomia de mecanismos estruturais, processuais e relacionais estabelecida 
por \cite{sengik_governanctecnologia_2022} fornece um arcabouço analítico sistemático para diagnosticar e avaliar a 
gestão pública. Essa abordagem possibilita verificar não apenas a existência de instâncias formais, 
mas também a maturidade das rotinas operacionais e a efetividade do relacionamento institucional 
indispensável para a geração de valor público por meio da tecnologia 
\cite{sengik_governanctecnologia_2022}.
}